\documentclass[10pt,a4paper]{article}
\usepackage{fontspec,amsmath,amssymb,graphicx}
\usepackage[margin=18mm]{geometry}
\setmainfont{Times New Roman}
\setlength{\parindent}{0pt}
\setlength{\parskip}{5pt}
\allowdisplaybreaks
\sloppy
\emergencystretch=2em
\begin{document}
\section*{BT\_6\_1\_1\_2\_01 — Lời giải bổ sung}
Dùng quy ước $f=a_0+\sum_{n\ge1}(a_n\cos k_nx+b_n\sin k_nx)$, $k_n=n\pi/L$, $a_0=(2L)^{-1}\int_{-L}^Lf$, $a_n=L^{-1}\int f\cos$, $b_n=L^{-1}\int f\sin$.
(i) $a_0=\sin(\alpha L)/(\alpha L)$, $a_n=2(-1)^n\alpha\sin(\alpha L)/[L(\alpha^2-k_n^2)]$, $b_n=0$. Ở $\alpha=0$ hoặc cộng hưởng, lấy giới hạn (hàm hằng hoặc một mode cos). Chuỗi hội tụ đều đến phép kéo dài liên tục chu kỳ $2L$.
(ii) $b_n=2(1-(-1)^n)/(n\pi)$, các $a_n=0$. Chuỗi hội tụ đến sign trong miền trừ điểm 0, về 0 tại 0 và các điểm nối $\pm L$; không hội tụ đều qua bước nhảy.
(iii) $a_0=\sinh(\alpha L)/(\alpha L)$,
\[a_n=\frac{2\alpha(-1)^n\sinh(\alpha L)}{L(\alpha^2+k_n^2)},\quad b_n=\frac{2k_n(-1)^{n+1}\sinh(\alpha L)}{L(\alpha^2+k_n^2)}.\]
Với $\alpha=0$, $f=1$. Với $\alpha\ne0$, chuỗi về $e^{\alpha x}$ trong $(-L,L)$, và $\cosh\alpha L$ tại mép; hội tụ đều trên compact nội, không trên toàn chu kỳ.
(iv) Phần ảo của $\log(1+re^{ix})$ cho $f=\sum_{n\ge1}(-1)^{n+1}r^n\sin nx/n$.
(v) Hai lần phần thực cho $f=2\sum_{n\ge1}(-1)^{n+1}r^n\cos nx/n$.
(vi) Nhân Poisson cho $f=(1-r^2)^{-1}[1+2\sum_{n\ge1}r^n\cos nx]$. Ở (iv)--(vi), $0<r<1$, chuỗi hội tụ tuyệt đối và đều theo tiêu chuẩn Weierstrass.\end{document}